\subsection{Location invariants on the trace model}
\label{sec:inv-semantics}
A location invariant maps every clock $x$ either to an upper bound
$M\in\mathbb R$, standing for the constraint $x\le M$, or to $+\infty$ (no
constraint).  We interpret automata with invariants on the extended words of
Section~\ref{sec:model}, as we do for automata without invariants.
Along a run $r$, the automaton stays in location $r(i)$ between the events
$i-1$ and $i$, and the invariant of $r(i)$ must hold during the whole stay,
that is, at the clock values $\CVal_\rho(i,x)-\tau$ for every delay
$\tau\in[0,\Delta_{i-1}]$, where $\Delta_{i-1}=t_i-t_{i-1}$ (and at the values $\CVal_\rho(0,x)$ for the
initial location).  This is the condition of Section~\ref{sec:model}, evaluated at
the event that ends the stay.  Since clocks increase while time elapses, an
invariant made of upper bounds holds during the stay if and only if it holds at
the event that ends it (lemma \texttt{inv\_\allowbreak{}during\_\allowbreak{}iff\_\allowbreak{}at\_\allowbreak{}event}).
For the two-sided invariants of Section~\ref{sec:forward}, a lower bound
$x\ge m$ must likewise hold during the whole stay, hence at its beginning.

\subsection{Invariant synthesis and backward propagation}
\label{sec:uppaal-invariants}\label{sec:main-invariants}
In the exported timed automaton, upper-bounded guards need invariants so that
time does not elapse past all useful outgoing transitions.  We compute location
invariants from outgoing upper bounds and propagate successor requirements
backward, using location invariants to constrain admissible time elapse as in
standard UPPAAL models~\cite{[BGD04]}.

\paragraph{Upper bound of a guard.}
For a conjunctive guard $g$ and a clock $x$, let $\mathit{ub}_x(g)$ be the
smallest bound $m$ of the constraints $x\le m$, $x<m$, and $x=m$ occurring in
$g$, and $+\infty$ if $g$ contains none.  Other constraints are ignored.  A
valuation satisfying $g$ satisfies $x\le\mathit{ub}_x(g)$.

\paragraph{Invariant synthesis.}
For a location $\ell$ with outgoing transitions
$\ell\xrightarrow{g_i,Z_i}\ell_i$, $i=1,\dots,m$, the synthesized invariant is
\[
  \mathit{Inv}(\ell)(x)=\max_{1\le i\le m}\mathit{ub}_x(g_i),
\]
with $\mathit{Inv}(\ell)(x)=+\infty$ when $\ell$ has no outgoing transition.
Thus a clock that some outgoing guard does not bound does not appear in the
invariant.  The bound is always non-strict, even when it comes from a strict
guard $x<m$: the invariant only needs to be implied by the guard of the next
transition, and strictness is kept in the guards themselves.

\begin{theorem}[Invariant synthesis]
\label{thm:inv-synthesis}
For every TBA $\mathcal A$, the automaton $\mathcal A$ with the synthesized
invariants accepts the same timed words as $\mathcal A$.
\end{theorem}

\begin{proof}
Adding invariants can only remove runs.  Conversely, let $r$ be an accepting
run on $\rho$.  At every event $i$, the transition taken from $r(i)$ is enabled,
so its guard holds at the values $\CVal_\rho(i,x)$, which are therefore at most
$\mathit{Inv}(r(i))(x)$.  By the lemma \texttt{inv\_\allowbreak{}during\_\allowbreak{}iff\_\allowbreak{}at\_\allowbreak{}event}, the invariant of $r(i)$ holds during
the whole stay.  (Rocq: \texttt{add\_\allowbreak{}invariants\_\allowbreak{}correct}.)
\end{proof}

\paragraph{Backward propagation.}
One propagation step strengthens every transition
$\ell\xrightarrow{g,Z}\ell'$ with the invariant of its target after the
resets:
\[
  g := g\wedge\mathit{Inv}(\ell')[Z:=0],
\]
where $\mathit{Inv}(\ell')[Z:=0]$ replaces every bound $x\le M$ on a reset clock
$x\in Z$ by the constant constraint $0\le M$, that is, by $\top$ if $M\ge0$ and
by $\bot$ otherwise.  Invariants are then synthesized again from the
strengthened guards.

\begin{theorem}[Backward propagation]
\label{thm:propagation}
For every TBA $\mathcal A$ and every $n\in\Nat$, the automaton obtained by
$n$ propagation steps accepts the same timed words as $\mathcal A$.
\end{theorem}

\begin{proof}
Strengthening guards can only remove runs.  Conversely, consider an accepting
run that takes $\ell\xrightarrow{g,Z}\ell'$ at event $i$ and then a transition
$\ell'\xrightarrow{g',Z'}\ell''$ at event $i+1$.  The second guard holds at event
$i+1$, so the values at that event are bounded by $\mathit{Inv}(\ell')$.  Clock
consistency gives $\CVal_\rho(i+1,x)=\Delta_i$ for $x\in Z$ and
$\CVal_\rho(i+1,x)=\CVal_\rho(i,x)+\Delta_i\ge\CVal_\rho(i,x)$ otherwise; hence
the strengthened guard $g\wedge\mathit{Inv}(\ell')[Z:=0]$ holds at event $i$.
The argument uses clock consistency, so the theorem is stated on timed words.
(Rocq: \texttt{propagate\_\allowbreak{}accepts} and
\texttt{propagate\_\allowbreak{}n\_\allowbreak{}accepts}.)
\end{proof}

Each round of the optimization pipeline of Section~\ref{sec:main-pipeline}
applies one propagation step.  The tool repeats the rounds until the
automaton is structurally equal to the previous one, as an OCaml value, or
until 50 rounds have been performed; the bound guarantees termination, and the
theorem holds for every number of steps, whether or not a fixed point is
reached. On the 29 completed optimizer inputs, the greatest observed number
of changed rounds is six (R6).
The invariant computation at one location is summarized in the article's
Figure~\ref{fig:main-inv-unified}.

\paragraph{Example.}
We illustrate the algorithm on the automaton of Fig.~\ref{fig:inv-example},
which contains a two-location cycle $\ell_1\to\ell_2\to\ell_1$ and a self-loop
at $\ell_2$.  It has one clock $x$, the initial location $\ell_0$, and the
transitions
\[
\begin{aligned}
&\ell_0 \xrightarrow{a,\;Z=\{x\}} \ell_1,\quad
 \ell_1 \xrightarrow{b,\;x\le 2,\;Z=\emptyset} \ell_2,\\
&\ell_2 \xrightarrow{c,\;\top,\;Z=\emptyset} \ell_1,\quad
 \ell_2 \xrightarrow{d,\;x\le 1,\;Z=\emptyset} \ell_2.
\end{aligned}
\]
The edge $\ell_1\xrightarrow{b}\ell_2$ requires $x\le2$, and the self-loop at
$\ell_2$ requires the tighter bound $x\le1$.

\begin{figure}[t]
\centering
\begin{tikzpicture}[node distance=21mm]
  \node[fmstate] (l0) {$\ell_0$};
  \node[fmstate,right=of l0] (l1) {$\ell_1$};
  \node[fmstate,right=of l1] (l2) {$\ell_2$};
  \draw[fmedge] ([xshift=-7mm]l0.west) -- (l0.west);
  \draw[fmedge] (l0) -- node[fmlabel,above] {$a\;/\;\mathsf{rst}(x)$} (l1);
  \draw[fmedge] (l1) to[bend left=18] node[fmlabel,above] {$x\le2,\;b$} (l2);
  \draw[fmedge] (l2) to[bend left=18] node[fmlabel,below] {$c$} (l1);
  \draw[fmedge] (l2) edge[loop above] node[fmlabel] {$x\le1,\;d$} (l2);
\end{tikzpicture}

\caption{Automaton used to illustrate invariant propagation, with one clock $x$
and actions $a$--$d$.  Edges are labeled ``guard, action / resets'' (the guard
is omitted when it is $\top$); only the edge $\ell_0\to\ell_1$ resets $x$.  The
incoming arrow marks the initial location; acceptance plays no role here.  The
computed fixed point is $\mathit{Inv}(\ell_1)=\mathit{Inv}(\ell_2)=(x\le2)$.}
\label{fig:inv-example}
\end{figure}

\begin{enumerate}
\item \emph{Synthesis.}  $\mathit{Inv}(\ell_1)(x)=\mathit{ub}_x(x\le2)=2$.  At
$\ell_2$, $\mathit{Inv}(\ell_2)(x)=\max(1,+\infty)=+\infty$, because the
transition toward $\ell_1$ does not bound $x$.  At $\ell_0$, the only outgoing
guard is $\top$, so $\mathit{Inv}(\ell_0)(x)=+\infty$.

\item \emph{Propagation.}  The transition
$\ell_2\xrightarrow{c,\;Z=\emptyset}\ell_1$ becomes guarded by
$\top\wedge\mathit{Inv}(\ell_1)=(x\le2)$.  The transition
$\ell_0\xrightarrow{a,\;Z=\{x\}}\ell_1$ receives $(0\le2)\equiv\top$ and is
unchanged.

\item \emph{Synthesis.}  At $\ell_2$, the outgoing guards are now $x\le1$ and
$x\le2$, so $\mathit{Inv}(\ell_2)(x)=\max(1,2)=2$.

\item \emph{Propagation.}  The transition $\ell_1\xrightarrow{b,\;x\le2}\ell_2$
receives $x\le2$, which it already contains, and the self-loop at $\ell_2$
receives $x\le2$, which is implied by $x\le1$.  Nothing changes.
\end{enumerate}
The fixed point is $\mathit{Inv}(\ell_0)=\top$ and
$\mathit{Inv}(\ell_1)=\mathit{Inv}(\ell_2)=(x\le2)$.  While the automaton
remains in the cycle $\ell_1\leftrightarrow\ell_2$, the clock $x$ cannot exceed
$2$, and at every reachable valuation at least one outgoing transition of the
current location is enabled.  The invariants do not remove every blocking
behavior: since $x$ is never reset in the cycle, time is blocked once $x=2$.
They make the upper bounds of the property automaton explicit, and
Theorems~\ref{thm:inv-synthesis} and~\ref{thm:propagation} guarantee that they
do not change its language.

\subsection{Forward propagation of lower bounds and simplification}
\label{sec:forward}\label{sec:main-forward}
\paragraph{Entry bounds.}
Dually, for a transition $t=\ell\xrightarrow{g,Z}\ell'$ and a clock $x$, the
\emph{entry bound} of $x$ along $t$ is $0$ if $x\in Z$ and otherwise the
largest bound $\mathit{lb}_x(g)$ of the constraints $x\ge m$, $x>m$, and $x=m$
of $g$, if any.  The entry bound of $x$ at a location $\ell'$ is the minimum of
the entry bounds of $x$ along the transitions entering $\ell'$; for the initial
location, the initial entry, with bound $0$, also counts.  If some entering
transition gives no bound on $x$, the clock $x$ has no entry bound at $\ell'$.
Since clocks are nonnegative and do not decrease while time elapses, every
clock with entry bound $m$ at $\ell'$ satisfies $x\ge m$ during every stay in
$\ell'$.

\paragraph{Forward propagation.}
The forward step adds $x\ge m$ to the guard of every transition leaving
$\ell'$, for every clock $x$ with entry bound $m$ at $\ell'$ (Rocq:
\texttt{forward\_\allowbreak{}accepts}).  The entry bounds can also be read as lower-bound
invariants: the two-sided automaton that carries the synthesized upper-bound
invariants and the entry bounds as lower-bound invariants accepts the same
timed words (\texttt{add\_\allowbreak{}two\_\allowbreak{}sided\_\allowbreak{}invariants\_\allowbreak{}accepts}).  UPPAAL
rejects lower-bound invariants, and their only role is to make infeasible
transitions visible: after forward and backward propagation, a transition whose
source location bounds a clock from below and whose target bounds it from above
carries both bounds in its guard.

\paragraph{Simplification.}
Three simplifications follow.  First, a transition whose guard contains
contradictory constraints is removed: a lower bound and an upper bound on the
same clock with $u<l$, or with $u=l$ and one of them strict, or an upper bound
$x\le u$ with $u<0$ or $x<u$ with $u\le0$.  Such a transition is never enabled (Rocq:
\texttt{remove\_\allowbreak{}contradictory\_\allowbreak{}accepts}).  Second, the transitions
leaving a non-initial location that has no entering transition
are removed: such a location never occurs on a run after position~$0$ (Rocq:
\texttt{remove\_\allowbreak{}unreachable\_\allowbreak{}accepts}).  The iteration of the pipeline
removes chains of such locations one step at a time.  A location whose only
entering transitions are self-loops keeps them, but it is unreachable from
the initial location; the tool counts and prints only the locations
reachable from the initial one and their transitions.  Third, a transition whose
action label is unsatisfiable, because it requires two different actions or an
action and its negation, is removed (Rocq:
\texttt{remove\_\allowbreak{}unsat\_\allowbreak{}accepts}).  Such labels arise because the
translator treats an action $a$ and the atom $\neg a$ of $\T(f)$ as independent
propositions; removing them also enables clock merging, since the resets of a
transition that can never be taken would otherwise prevent it.

\subsection{Dead resets and merging of synchronously reset clocks}
\label{sec:dead-merge}\label{sec:main-dead-merge}
\paragraph{Dead resets.}
A clock $x$ is \emph{live} at a location $\ell$ if some transition that tests
$x$ in its guard can be taken from $\ell$, possibly after other transitions.
The live locations are computed as a backward closure, along all
transitions, from the sources of the transitions that test $x$; the other
locations are \emph{dead} for $x$.  Since the correctness of the pass does
not depend on how this closure is computed, the pass checks the two
properties that the proof uses, namely that no transition leads from a dead
location to a live one and that none tests $x$ from a dead location, and then removes $x$ from the reset sets of the
transitions that enter a dead location.  After such a transition, the value of
$x$ is never read again, so whether it is reset does not matter.  The proof
builds, from every accepted extended word, an extended word that differs only
in the reset decisions of $x$, and conversely (Rocq:
\texttt{remove\_\allowbreak{}dead\_\allowbreak{}resets\_\allowbreak{}clock\_\allowbreak{}accepts}).  The pass is applied
to every clock (\texttt{remove\_\allowbreak{}dead\_\allowbreak{}resets\_\allowbreak{}accepts}).

\paragraph{Synchronous clocks.}
The verified translation shares the clock of syntactically identical timed
subformulas.  Clocks of \emph{different} subformulas can also be merged when
they are reset synchronously.  Two clocks $x\neq y$ are \emph{mergeable} if
\begin{enumerate}[label=(\roman*)]
  \item every transition resets $x$ if and only if it resets $y$;
  \item no transition enters the initial location; and
  \item every transition leaving the initial location resets $x$ and does not
        test it.
\end{enumerate}
Then, after the first event, $x$ and $y$ always have the same value, and every
test of $x$ can be replaced by the same test on $y$ (Rocq:
\texttt{merge\_\allowbreak{}accepts}).  The condition is checked by a Boolean function
proved to imply it, and the pass tries all pairs of reset clocks
(\texttt{merge\_\allowbreak{}all\_\allowbreak{}accepts}).  The merged clock $x$ is no longer tested,
so its remaining resets are removed by the dead-reset pass in the next round.
This verified merging is restricted to synchronously reset clocks; merging
clocks whose values coincide for other reasons is not covered.

\subsection{Normalization of guards}
\label{sec:normalize}
Backward and forward propagation add constraints to guards.  Every time a
guard is strengthened, it is \emph{tightened}: a constraint implied by another
constraint of the same guard on the same clock is dropped, so that a guard
keeps at most one lower bound and one upper bound per clock (an equality
counting as both); for instance, $x\le5\land x<2$ becomes $x<2$.  The
implication test compares bounds and strictness, and is proved sound (Rocq:
\texttt{tighten\_\allowbreak{}holds}).  Normalization tightens every guard and also
removes the lower bounds $x\ge m$ with $m\le0$ and $x>m$ with $m<0$, which every
clock value satisfies; it is needed after the merging of clocks, which renames
constraints.  The removal of trivial lower bounds relies on the nonnegativity
of clock values, so the theorem is stated on timed words (Rocq:
\texttt{normalize\_\allowbreak{}accepts}).

\subsection{Merging of transitions and of locations}
\label{sec:merge-ts}\label{sec:main-merge-ts}
\paragraph{Transitions.}
Two transitions with the same source, target, and resets are merged when the
union of their enabling conditions is again a conjunction.  A transition is
removed when another transition with the same source, target, and resets is
\emph{weaker}, that is, when its label and guard imply those of the other
transition, constraint by constraint (subsumption).  Two transitions
whose labels, or guards, differ only by a literal and its negation, such as
$p$ and $\neg p$, $x<2$ and $x\ge2$, or $x\le5$ and $x>5$, are replaced by one
transition without that literal (resolution).  The pass preserves the language
because every enabled transition has an enabled counterpart with the same
source and target, and conversely (Rocq: \texttt{merge\_\allowbreak{}trans\_\allowbreak{}accepts}).
The other unions are disjunctive; they are formed by the export
(Section~\ref{subsec:split-construction}).

\paragraph{Locations.}
A location $\ell'$ is merged into a location $\ell$ when both are accepting or
both are not, and when every transition leaving $\ell'$ has a twin leaving
$\ell$, and conversely: same label, guard, and resets, compared as sets, and
same target once $\ell'$ is replaced by $\ell$.  The transitions entering
$\ell'$ are redirected to $\ell$, and those leaving $\ell'$ are removed.  A run
of the original automaton becomes a run of the merged one by replacing $\ell'$
by $\ell$; conversely, a run of the merged automaton is followed in the
original one, taking in $\ell'$ the twin of the transition taken in $\ell$
(Rocq: \texttt{merge\_\allowbreak{}states\_\allowbreak{}all\_\allowbreak{}accepts}; ``state'' in a Rocq name
means location).  The condition on the B\"uchi marking is necessary: merging
an accepting location with a non-accepting twin would accept runs that visit only the latter infinitely
often.
